\section{Inductive Proofs and Cyclic proofs}

This section defines cyclic proof systems and inductive proof systems.

\subsection{Intuitionistic Martin-L\"of's Inductive Definition System $\LJID$}

In this section,
we define
an intuitionistic Martin-L\"of's inductive definition system, called $\LJID$,
which is an intuitionistic version of $\LKID$ given in \cite{Brotherston11}
and formalizes intuitionistic first-order logic with inductive definitions
in sequent calculus.

The language of $\LJID$ is a first-order language with
production rules for inductive predicate symbols.
The predicate symbols are classified into ordinary predicates symbols $Q,Q_1,Q_2,\ldots$
and inductive predicate symbols $P,P_1,P_2,\ldots$.
We assume the first order terms $t,u, \ldots$.
We call an atomic formula an {\em inductive atomic formula}
when its predicate symbol is an inductive predicate symbol.

For a formula $F$ and a variable $x$,
we will call the expression $\lambda x.F$ an {\em abstract}.
We will sometimes define an abstract $G$ as $\lambda x.F$
by defining $G(x)$ as the formula $F$.
For an abstract $\lambda x.F$ and a term $t$,
we write $(\lambda x.F)t$ for $F[x:=t]$.
We call a predicate symbol or an abstract a {\em predicate}.

For simplicity,
we sometimes write $Fxy$ for $F(x,y)$ when $F$ is a predicate,
We write $e[x:=t]$ for the expression obtained from $e$ by replacing
every free occurrence of the variable $x$ by $t$.
For simplicity,
We also write $e[\Vec x:=\Vec t]$ for simultaneous multiple substitution.
We also write $\theta$ for the substitution $[\Vec x:=\Vec t]$.
We write $\FV(\Vec e)$ for the set of free variables in the expressions $\Vec e$.

A {\em production rule} is defined as
\[
\infer{P(\Vec t)}{
F_1 \ \ldots \ F_n
}
\]
where $F_1,\ldots,F_n$ are atomic formulas and
$P$ is an inductive predicate symbol.
The inductive predicate symbol $P_i$ is {\em dependent} on 
the inductive predicate symbol $P_j$,
if
$i=j$ or
there are some $k$ and some production rule
such that 
the conclusion of the production rule contains $P_i$ and
the assumptions of the production rule contain $P_k$
and $P_k$ is dependent on $P_j$.
The inductive predicate symbols $P_i,P_j$ are {\em mutually dependent}
if $P_i$ is dependent on $P_j$ and $P_j$ is dependent on $P_i$.

These production rules mean that 
the inductive predicate denotes 
the least fixed point defined by these production rules.
%
For example, 
the production rules of the inductive predicate symbol $N$ are given by
\[
\infer{N0}{}
\qquad
\infer{Ns x}{Nx}
\]
These production rules mean that 
$N$ denotes the smallest set that contains 0 and closed under $s$,
namely the set of natural numbers.

$\LJID$ is based on a sequent calculus.
We use a sequent $\Gamma \prove \Delta$ where $\Gamma$ and $\Delta$
are sequences of formulas and
the length of $\Delta$ is $\le 1$.

The inference rules of $\LJID$ are those of $\LJ$ (given in the appendix)
and the following $(P\ R)$ and $(\Ind\ P_j)$.

{\prooflineskip
\[
\infer[(P\ R)]{\Gamma \prove P\Vec t}{
\Gamma \prove Q_{1}\Vec u_{1} & \ldots & \Gamma \prove Q_{n}\Vec u_{n}
& \Gamma \prove P_1\Vec t_{1} & \ldots & \Gamma \prove P_{m}\Vec t_{m}
}
\\
\infer[(\Ind\ P_j)]{\Gamma,P_j\Vec u \prove \Delta}{
\hbox{minor premises}
&
\Gamma,F_j\Vec u \prove \Delta
}
\]
}
For $(P\ R)$ we assume the production rule
\[
\infer{P\Vec t}{Q_1\Vec u_1 & \ldots & Q_n\Vec u_n & P_1\Vec t_1 & \ldots & P_m\Vec t_m}
\]
In the figure, 
for $(\Ind\ P_j)$, we also assume 
a predicate $F_i$ for each $P_i$,
and
the minor premises are defined as 
\[
\Gamma, Q_{i1}\Vec u_{i1}, \ldots, Q_{in_i}\Vec u_{in_i}, 
F_{i1}\Vec t_{i1}, \ldots, F_{im_i}\Vec t_{im_i}
\prove F_i\Vec t_i
\]
for each production rule
\[
\infer{P_i\Vec t_i}{Q_{i1}\Vec u_{i1} & \ldots & Q_{in_i}\Vec u_{in_i} & P_{i1}\Vec t_{i1} & \ldots & P_{im_i}\Vec t_{im_i}}
\]
where $P_i$ is mutually dependent on $P_j$.
An example of these rules is in the appendix.

The rules $(\neg L)$, $(\neg R)$,
$(\land L)$, $(\land R)$,
$(\lor L)$, $(\lor R1)$, $(\lor R2)$, 
$(\imp L)$, $(\imp R)$,
$(\forall L)$, $(\forall R)$,
$(\exists L)$, and $(\exists R)$
are called {\em logical rules}, and
the introduced formulas in these rules and $(P\ R)$
are called {\em main formulas},
and the corresponding formulas in the assumptions are called {\em auxiliary
formulas}.

\subsection{Cyclic Proof System $\CLJIDomega$}

This section defines
a cyclic proof system $\CLJIDomega$ for intuitionistic first-order logic
with inductive definitions.

$\CLJIDomega$ is an intuitionistic version of $\CLKIDomega$ 
defined in \cite{Brotherston11}.
The inference rules of $\CLJIDomega$ are obtained from
those of $\LJID$ by replacing $(\Ind\ P_j)$ by
\[
\infer[(\Case P)]{\Gamma, P\Vec u \prove \Delta}{
\hbox{case distinctions}
}
\]
where the case distinctions are
\[
\Gamma,\Vec u=\Vec t,
Q_1\Vec u_1, \ldots, Q_n\Vec u_n, P_1\Vec t_1, \ldots, P_m\Vec t_m
\prove \Delta
\]
for each production rule
\[
\infer{P\Vec t}{Q_1\Vec u_1 & \ldots & Q_n\Vec u_n & P_1\Vec t_1 & \ldots & P_m\Vec t_m}
\]
and
$\FV(\Gamma,P_j\Vec u, \Delta)
\cap \FV(P\Vec t,Q_1\Vec u_1, \ldots, Q_n\Vec u_n, P_1\Vec t_1, \ldots, P_m\Vec t_m) = \emptyset$.

For the rule $(\Case P)$,
the introduced formula 
is called a {\em main formula},
and the corresponding formulas in the assumptions are called {\em auxiliary
formulas}.

\begin{Def}[Preproof]\rm
A {\em preproof} in $\CLJIDomega$ is defined as 
a finite figure generated by combination of inference rules,
where 
we allow an open assumption (called a {\em bud}) and
requiring for each bud
some corresponding sequent of the same form (called a {\em companion})
that is not a bud.
\end{Def}
We call the correspondence from a bud to a companion {\em bud-companion relation}.
%
An {\em infinite unfolding} of a preproof is defined as
the infinite figure obtained by 
possibly infinitely many times replacing each bud by
the subproof of the companion of the bud 
to reach some possibly infinite figure without any buds.

\begin{Def}\rm
A {\em path} in a preproof or an infinite unfolding of a preproof
is defined as a sequence of sequents obtained by going
upwardly from some node to another node in the tree
of a proof.
An {\em infinite path} in an infinite unfolding of a preproof
is defined similarly.
We define a tail of a path as the sequence of sequents above
some sequent in the path.

For a path and an infinite path,
we define
a {\em trace}
as
a sequence of inductive atomic formulas 
in the antecedents of sequents in the path, 
constructed by
chasing a corresponding inductive atomic formula upwardly.
More precisely,
we chase an auxiliary inductive atomic formula when 
the inductive atomic formula is the main formula of $(\Case P)$,
and
we chase a corresponding inductive atomic formula for
the rules $(\Subst)$ and $(=L)$ and $(\Wk\ L)$ and $(\Wk\ R)$ and 
$(\Cont)$ and $(\Exch)$.
We chase the corresponding same inductive atomic formula upwardly in the other cases.
Namely,
in $(\Case P)$ we chase $P_i\Vec t_i$ after $P\Vec t$,
in $(\Subst)$ we chase $F$ after $F\theta$,
and
in $(=L)$ we chase $F[x:=t]$ after $F[x:=u]$.

We call a trace {\em progressing} 
if it passes some main formula of a case rule.
We call a trace {\em infinitely progressing} 
if it passes main formulas of case rules infinitely many times.

The {\em global trace condition} \cite{Brotherston11,Brotherston06} is 
the condition that
for every infinite path in the infinite unfolding of a given cyclic proof,
there is some infinitely progressing trace in some tail of the infinite path.
\end{Def}

We define cyclic proofs.
An example of cyclic proofs is given in the appendix.

\begin{Def}[Cyclic Proof]\rm
A {\em proof} (called a {\em cyclic proof}) in $\CLJIDomega$ is defined 
as a preproof that satisfies the global trace condition.
\end{Def}

